(1) 从给定状态方程求 $(\partial p/\partial T)_V=R/(V-b)$，故 $T(\partial p/\partial T)_V-p=T R/(V-b)-[RT/(V-b)-a/V^2]=a/V^2$。积分得 $U(T,V)=C_V T-a/V+U_{\rm ref}$；参数 $b$ 虽影响压强，当前模型的内能体积依赖只含吸引项 $-a/V$。
(2) 对绝热自由膨胀，$Q=0,W=0$，第一定律给 $\Delta U=0$。于是 $C_V(T_f-T_i)=a(1/V_f-1/V_i)<0$，即 $T_f=T_i+(a/C_V)(1/V_f-1/V_i)$。必须满足 $T_f>0$、$V_i,V_f>b$ 且末态仍处于所用单一气相区；否则范德瓦耳斯单相模型与恒定热容假设需重新评估。
(3) 恒温准静态路径的功为 $W_T=\int_{V_i}^{V_f}[RT/(V-b)-a/V^2]\,dV=RT\ln[(V_f-b)/(V_i-b)]+a(1/V_f-1/V_i)$。其内能变化 $\Delta U_T=a(1/V_i-1/V_f)>0$，故第一定律给 $Q_T=\Delta U_T+W_T=RT\ln[(V_f-b)/(V_i-b)]$。温度保持不变需要向系统输入热，完全不同于绝热自由膨胀。
(4) 当 $a\to0$，内能只依赖 $T$，自由膨胀温度变化趋零；当 $V_f\to V_i$，温度变化、积分功与热量均连续趋零。气体膨胀时平均分子吸引势能 $-a/V$ 上升，若容器隔热且没有外功，只能通过降低分子热运动内能补偿，因此末温下降。
